\section{Verification coordinates for common reversible responses}\label{sec:supp-r72}
In the compact-group argument the map from hidden rows to response functions is defined only on the reachable row span $R$. A whole realizing simplex need not lie in that span. Thus the admissible polytope is the image of $R\cap\Delta_M$, not an unspecified positive cone on the minimal physical support. The support-count estimate uses the fact that two distinct extreme points of the section cannot have exactly the same positive coordinates. It is only a size certificate for a supplied realization, not a constructive upper on minimal stochastic memory.

The quantum specialization removes the common kernel of all commanded terminal tests. Each unitary conjugation restricts to an isometry on the finite-dimensional reachable Hermitian span. Forward invariance and finite dimension imply backward invariance, both for that span and the invisible kernel. The quotient therefore has a common invariant inner product. No controllability or informational completeness of the terminal measurement is assumed; a constant effect has a trivial observable group even if the full physical group is infinite.

For the two reflections the physical chronological convention is $\Phi_{a_n}\cdots\Phi_{a_1}\rho$, while the observer row is $e_1^{\mathsf T}Q_{a_1}\cdots Q_{a_n}$. The intertwining equations use row vertices and the transpose of the physical Bloch map. This is why the word $(01)^n$ corresponds to powers of $Q_0Q_1$ and to the rotation produced by the physical map $S_1S_0$. The explicit diagnostic uses both equations and a separate raw density-matrix evolution. Reversing only one convention would break the mixed-word identity.

\section{Partition synthesis, anchors and physical interpretation}\label{sec:supp-r72-partition}
The initialization row in the periodic construction chooses its recurrent block only after the first actual transition. At time zero it uses a separately assigned decoder $b_0$. Entering phase one, rather than phase zero, ensures the first observed response is correct. The lower's extra label is not an arbitrary implementation preference: two reachable closed recurrent classes cannot contain the same prescribed initial state. A model permitting an initially mixed state is a different resource convention and can remove that label.

In the decay-normalized lower, the full nonnegative core need not be irreducible, row stochastic or similar by one strictly positive global Perron vector to a stochastic matrix. Its powers are bounded because every active entry is bounded by a fixed shifted target-output path. Perron scaling is performed only on a critical irreducible diagonal block. The no-critical-chain lemma forbids recovering an initialization label by sending a critical class into another required critical class. The proof counts nonnegative paths, so cancellation cannot conceal the resulting quadratic cumulative growth.

The periodic response has a prepared block-qubit realization and a fixed terminal effect. The established flag embedding realizes the associated exact-response requirement using one fixed squared loss and three actual audit outcomes at the original first-herald budget. Internal damping, physical leakage and an observed terminal event are different operations. The new finite-memory laws concern exact replication of the terminal probability, not predicting the realized hidden bit or the distinct squared-sign decision task. They do not give the protected system to the observer.

\section{Independent diagnostics for the present revision}\label{sec:supp-r72-diagnostics}
The new program compares subset dynamic programming with an independent enumeration of all set partitions on finite input lists. Rational and algebraic examples check row normalization, bounded decoders, the initialization phase, exact damped and leaked coefficients and random-horizon resolvents. A positive substochastic two-block example shows exactly how linking critical classes yields linearly unbounded rescaled powers. These checks do not replace the class-decomposition lower proofs.

For the controlled theorem, a shared three-vertex algebraic construction verifies both reflection matrices and their intertwining identities exactly. Every tested mixed word is independently evolved through the raw qubit channels. Additional numeric polygon constructions check longer binary words, both physical maps, their pairwise product and phased periodic mixtures. They establish finite implementation consistency only. A finite angular or continued-fraction sample is never used to infer an irrationality exponent or uniform asymptotic class.
